%%%%%%%%%%%%%%%%%%%%%%%%%%%%%%%%%%%%%%%%%%%%%%%%%%%%%%%%%%%%%%%%%%%%%%%%%%%%%%%
\subsubsection{Domain 10.9.4: Sustainability and Climate Behavior}
\label{subsubsec:sustainability}
%%%%%%%%%%%%%%%%%%%%%%%%%%%%%%%%%%%%%%%%%%%%%%%%%%%%%%%%%%%%%%%%%%%%%%%%%%%%%%%

\paragraph{The Challenge.}
Climate change requires individuals to adopt pro-environmental behaviors---
reducing consumption, changing transportation, dietary shifts---that impose
personal costs for collective benefits. The INU-KNU conflict is stark:
individual sacrifice for global good.

\paragraph{Context Profile ($\Psi$).}

\begin{table}[htbp]
\centering
\caption{Sustainability Context Factors}
\label{tab:sustainability-psi}
\renewcommand{\arraystretch}{1.3}
\begin{tabular}{@{}lll@{}}
\toprule
\textbf{$\Psi$-Dimension} & \textbf{Level} & \textbf{Effect} \\
\midrule
$\Psi_S$ (Social trust) & Critical & Low: ``Others won't cooperate'' \\
$\Psi_T$ (Temporal) & Problematic & Climate = long-term, heavily discounted \\
$\Psi_K$ (Information) & Variable & Climate literacy varies widely \\
$\Psi_I$ (Institutional) & Mixed & Policies can enable or block action \\
\bottomrule
\end{tabular}
\end{table}

\paragraph{The Core Conflict: INU vs. KNU/Eco.}
Pro-environmental behavior creates a direct category conflict:

\begin{align}
\Delta U^{INU} &< 0 \quad \text{(personal cost: money, convenience, comfort)} \\
\Delta U^{KNU} &> 0 \quad \text{(collective benefit: shared future)} \\
\Delta U^{Eco} &> 0 \quad \text{(environmental quality)}
\end{align}

\paragraph{The Scale Mismatch Problem.}
Individual action has negligible global impact:

\begin{equation}
\text{My contribution to climate solution} \approx \frac{1}{8 \text{ billion}} \approx 0
\end{equation}

So pure consequentialist INU reasoning yields: ``My action doesn't matter.''

\paragraph{The Temporal Mismatch Problem.}
Climate benefits are far-future and uncertain:

\begin{align}
\text{Cost:} &\quad U^{INU}_{F,0}, U^{INU}_{E,0} \quad \text{(immediate, certain)} \\
\text{Benefit:} &\quad U^{KNU}_{Eco,3}, U^{KNU}_{P,3} \quad \text{(distant, uncertain)}
\end{align}

With heavy discounting and uncertainty:
\begin{equation}
\text{Present value of climate benefit} \approx \delta^3 \cdot p(\text{success}) \cdot \text{share} \approx 0
\end{equation}

\paragraph{Why Do People Act Pro-Environmentally?}
If INU analysis predicts inaction, why do some people act? The answer
lies in KNU and especially IDN:

\begin{table}[htbp]
\centering
\caption{Motivations for Pro-Environmental Behavior}
\label{tab:env-motivations}
\renewcommand{\arraystretch}{1.3}
\begin{tabular}{@{}llp{5cm}@{}}
\toprule
\textbf{Category} & \textbf{Component} & \textbf{Mechanism} \\
\midrule
KNU & $U^{KNU}_{Eco}$ & ``For my grandchildren'' \\
KNU & $U^{KNU}_{S}$ & ``My community does this'' \\
IDN & $U^{IDN}_E$ & ``I'm an environmentalist'' \\
IDN & $U^{IDN}_S$ & ``People like me care about climate'' \\
INU & $U^{INU}_E$ & ``I feel good when I act green'' \\
\bottomrule
\end{tabular}
\end{table}

\paragraph{The Identity Channel.}
Environmental identity ($U^{IDN}$) is often the primary driver:

\begin{equation}
r^{IDN}_{Eco} = \text{``An environmentalist has carbon footprint } < X\text{''}
\end{equation}

For someone with strong environmental identity:
\begin{itemize}
\item Flying creates $P^{IDN}$ (identity pain)
\item Cycling creates $G^{IDN}$ (identity confirmation)
\item These effects are \emph{immediate} (not discounted)
\end{itemize}

\paragraph{Welfare Calculation: Reducing Meat Consumption.}
Consider a decision to reduce meat consumption:

\begin{table}[htbp]
\centering
\caption{Meat Reduction: Welfare Analysis}
\label{tab:meat-welfare}
\small
\renewcommand{\arraystretch}{1.2}
\begin{tabular}{@{}lrrr@{}}
\toprule
\textbf{Component} & \textbf{$\Delta U$} & \textbf{Weight} & \textbf{Contribution} \\
\midrule
\multicolumn{4}{l}{\textit{INU}} \\
$U^{INU}_{E,0}$ (taste pleasure) & --0.30 & 0.40 & --0.12 \\
$U^{INU}_{F,0}$ (cost savings) & +0.10 & 0.40 & +0.04 \\
$U^{INU}_{P,2}$ (health) & +0.15 & 0.40 & +0.06 \\
\midrule
\multicolumn{4}{l}{\textit{KNU}} \\
$U^{KNU}_{S,0}$ (social friction) & --0.20 & 0.25 & --0.05 \\
$U^{KNU}_{Eco,3}$ (climate benefit) & +0.01 & 0.25 & +0.003 \\
\midrule
\multicolumn{4}{l}{\textit{IDN --- Environmental Identity HIGH}} \\
$U^{IDN}_{E,0}$ (identity confirmation) & +0.50 & 0.35 & +0.18 \\
\midrule
\multicolumn{3}{l}{\textbf{Total $\Delta Q$ (environmental identity)}} & \textbf{+0.11} \\
\midrule
\multicolumn{4}{l}{\textit{IDN --- Environmental Identity LOW}} \\
$U^{IDN}_{E,0}$ (no identity effect) & 0 & 0.35 & 0 \\
\midrule
\multicolumn{3}{l}{\textbf{Total $\Delta Q$ (no environmental identity)}} & \textbf{--0.07} \\
\bottomrule
\end{tabular}
\end{table}

\textbf{Key finding:} Environmental identity determines behavior.
With strong IDN, meat reduction improves welfare.
Without IDN, it reduces welfare.

\paragraph{Empirical Evidence: The Meat Paradox.}
Recent cross-country research \citep{giner2025sustainable} provides compelling evidence for this identity-behavior dynamic. Using latent class analysis across 9 OECD countries (n=8,261), Giner et al. (2025) identify four distinct consumer segments:
\begin{itemize}
\item \textbf{Class 1 (41\%):} Low environmental concern, high meat consumption
\item \textbf{Class 2 (21\%):} Moderate concern, moderate consumption
\item \textbf{Class 3 (26\%):} High environmental concern, \emph{high} meat consumption --- the ``Meat Paradox''
\item \textbf{Class 4 (12\%):} High concern, low consumption (identity-aligned)
\end{itemize}

Class 3 exemplifies cognitive dissonance ($CDI = 0.26$): awareness ($U^{KNU}_{Eco} > 0$) without identity integration ($U^{IDN}_{E} \approx 0$). These individuals acknowledge climate concerns but lack the identity commitment that would make behavioral change welfare-improving. Notably, policy support varies dramatically: meat tax support ranges from 16\% (Class 1) to 36\% (Class 4), suggesting identity-based segmentation for intervention design.

\paragraph{The Social Tipping Point.}
KNU effects depend on social norms:

\begin{equation}
U^{KNU}_{S} = f(\text{share of others doing same})
\end{equation}

\begin{itemize}
\item If few people are vegetarian: Social friction ($U^{KNU}_S < 0$)
\item If many are vegetarian: Social belonging ($U^{KNU}_S > 0$)
\end{itemize}

This creates a \emph{tipping point}---once enough people adopt, 
the social cost inverts to social benefit.

\paragraph{Intervention Strategies.}

\begin{table}[htbp]
\centering
\caption{Sustainability Interventions}
\label{tab:sustainability-interventions}
\renewcommand{\arraystretch}{1.3}
\begin{tabular}{@{}lll@{}}
\toprule
\textbf{Target} & \textbf{Strategy} & \textbf{Example} \\
\midrule
$U^{IDN}$ & Build environmental identity & ``Climate-conscious consumer'' labels \\
$U^{KNU}$ & Social proof & ``73\% of neighbors reduced energy'' \\
$U^{INU}_{F}$ & Align financial incentives & Carbon tax + dividend \\
$a_{KNU}$ & Activate collective awareness & ``For our grandchildren'' framing \\
$\gamma_{INU,KNU}$ & Create positive complementarity & Community solar (individual + collective gain) \\
\bottomrule
\end{tabular}
\end{table}

\paragraph{The $\Gamma$-Matrix Challenge.}
Recall from Section~\ref{subsec:context-modulation}:
\begin{align}
\gamma_{Eco,M} &= -0.21 \quad \text{(markets harm environment)} \\
\gamma_{Eco,E} &= -0.19 \quad \text{(development harms environment)}
\end{align}

These structural trade-offs mean that standard economic growth 
\emph{systematically} undermines ecological welfare. Policy must 
actively counteract these forces.

\paragraph{Case Example: Social Norm Messaging for Energy Conservation.}
Opower's home energy reports show:
\begin{itemize}
\item Your energy use vs. neighbors (social comparison → $r^{INU}$ shift)
\item Smiley faces for below-average use (immediate $G^{INU}_E$)
\item ``Efficient homes like yours'' framing ($U^{IDN}$)
\end{itemize}

Result: 2--3\% energy reduction with no financial incentive---pure
KNU and IDN activation.

\paragraph{Key Insight.}
\begin{tcolorbox}[colback=yellow!10!white, colframe=yellow!50!black]
\textbf{Sustainability Insight:}
Pro-environmental behavior cannot be driven by INU alone:
\begin{enumerate}
\item Individual impact is negligible (scale mismatch)
\item Benefits are distant and uncertain (temporal mismatch)
\item INU costs are immediate and certain
\end{enumerate}
Successful sustainability interventions must:
\begin{itemize}
\item Build environmental identity ($U^{IDN}$)
\item Activate collective concern ($U^{KNU}$)
\item Create positive INU-KNU complementarity (win-win designs)
\item Use social proof to shift norms (tipping point)
\end{itemize}
The $\Gamma$-matrix trade-offs ($\gamma_{Eco,M} < 0$, $\gamma_{Eco,E} < 0$) 
mean this requires active policy intervention, not market forces alone.
\end{tcolorbox}

%%%%%%%%%%%%%%%%%%%%%%%%%%%%%%%%%%%%%%%%%%%%%%%%%%%%%%%%%%%%%%%%%%%%%%%%%%%%%%%
% END OF 10.9.4
%%%%%%%%%%%%%%%%%%%%%%%%%%%%%%%%%%%%%%%%%%%%%%%%%%%%%%%%%%%%%%%%%%%%%%%%%%%%%%%
